\section{Introduction}
\label{sec:introduction}

Retrieval-augmented generation (RAG) over documents with rich visual layout -- web pages, tables, figures, infoboxes -- has long relied on text extraction: parse the HTML or run OCR, chunk the resulting text, and retrieve chunks with a dense text retriever. This pipeline discards exactly the information that layout encodes: which numbers belong to which row of a table, which caption belongs to which figure, which heading scopes which paragraph. A parallel line of work asks whether a vision-language model (VLM) can skip text extraction altogether and read the document as an image. PixelRAG~\citep{pixelrag2026} shows that retrieving directly over rendered page screenshots, rather than extracted text, can match or beat text-based retrieval on long, layout-heavy documents.

PixelRAG's tiling, however, is a fixed pixel grid applied uniformly to the rendered page, independent of what is actually on it. Independent coverage of the method notes the direct consequence: the grid ``slices pages by fixed pixel height, meaning a table or paragraph can get cut in half mid-tile with no awareness of content boundaries''\todo{cite VentureBeat coverage precisely, or drop this quote if a primary source is preferred}. A retriever built on top of such tiles has to reconstruct content boundaries that the renderer already knew and threw away.

MAGE-RAG~\citep{magerag2026} addresses a related but distinct problem: given a long, already-segmented document, how should an agent decide \emph{which} passages to actually read under a bounded reasoning budget? It builds a multigranular evidence graph over page- and element-level nodes with several explicit relations (containment, reading order, layout adjacency, section hierarchy, semantic neighbors) and an online controller (its Algorithm~1) that walks the graph, activating and pruning nodes, so that only a budgeted subset of the document is ever read by the downstream model. MAGE-RAG assumes the document is already parsed into this structure; it does not address how to obtain content-aligned tiles from a rendered page in the first place.

These two systems are complementary: PixelRAG shows that pixel tiles read by a VLM are a viable unit of retrieval for visually rich documents, but tiles them blindly; MAGE-RAG shows that a graph-structured, budget-aware controller over document elements outperforms retrieving a static top-$k$, but assumes the elements already exist. This paper combines them by using the one artifact PixelRAG discards and MAGE-RAG never had to construct in the first place: the DOM. When the input is a rendered web page rather than a scanned PDF, the browser's own layout engine already knows the exact bounding box of every paragraph, heading, table, and image before a single pixel is cropped.

\paragraph{Contributions.} We introduce \textbf{DOM-Graph RAG}, an open-source pipeline (\S\ref{sec:method}) that:
\begin{enumerate}
\item replaces PixelRAG's blind pixel grid with \emph{DOM-guided adaptive tiling}: elements are cropped along their own bounding boxes, undersized elements (e.g.\ short list items) are merged with their vertical neighbors, oversized elements (e.g.\ a very tall table) are split into sub-tiles, heading elements are never emitted alone but always paired with the content they introduce, and elements that nest inside an already-kept element's bounding box (e.g.\ an \code{<img>} inside a \code{<figure>}) are pruned before tiling rather than tiled twice (\S\ref{sec:tiling});
\item organizes these tiles into a \emph{multigranular graph} in the style of MAGE-RAG -- page and element nodes connected by containment, spatial reading-order, layout-adjacency, and section-hierarchy relations -- and extends it across documents via hyperlink edges into a small corpus graph, rather than a single isolated page; a page too long to tile as one unit of work is itself \emph{split into section-level sub-pages}, linked by a new \code{continues} relation and merged with the same node-namespacing machinery as the corpus graph, so that document length is handled by growing the graph rather than by excluding long documents (\S\ref{sec:graph});
\item fine-tunes a general-purpose VLM into a page-element reader/retriever via \emph{contrastive LoRA training} on synthetically generated (tile, question) pairs, using hard negatives mined by lexical similarity and, optionally, by a VLM judge that inspects each candidate tile directly to catch non-lexical false negatives (a paraphrase, or an answer conveyed visually rather than as text), with adapters on both the language decoder and the vision tower; the same tile set is also mirrored at reduced resolution via Lanczos downsampling, so that training can be run and compared at full and reduced resolution as a direct, controlled test of PixelRAG's resolution/token-cost tradeoff claim (\S\ref{sec:contrastive});
\item implements an \emph{online evidence controller} that evolves each element node through an explicit four-state machine (\state{inactive}, \state{active}, \state{opened}, \state{pruned}) via a best-first search under a node-opening budget, generalizing MAGE-RAG's Algorithm~1 to graphs built from live DOM structure rather than pre-parsed PDF layout (\S\ref{sec:controller}).
\end{enumerate}
The full pipeline -- rendering, tiling, graph construction (including oversized-page splitting and multi-page corpus composition), contrastive fine-tuning (including the resolution ablation), and evidence control -- is implemented and exercised end-to-end on live Wikipedia pages in an accompanying notebook and a standalone batch script, with dependency-free unit tests covering the geometric, graph, mining, and image-processing logic. \S\ref{sec:experiments} lays out the intended evaluation against MAGE-RAG's published LongDocURL and MMLongBench-Doc baselines; we report this as work in progress rather than fabricate numbers ahead of running it.
